\section{Modelos de lenguaje y modelos visión--lenguaje}
\label{sec:marco_llm}

Generar un reporte a partir de una imagen requiere tres piezas: un modelo capaz de
producir texto, una forma de adaptarlo a la dermatología sin reentrenar todos sus parámetros y una forma de
conectarlo con la imagen. La primera de esas piezas es el modelo de lenguaje.

\subsection{Grandes modelos de lenguaje}
\label{sec:marco_llm_grandes}

Un \textbf{modelo de lenguaje} asigna a cada token del vocabulario la probabilidad de continuar un texto dado.
Por ejemplo, después de «La lesión presenta una red de pigmento», una continuación como «atípica» recibe más
probabilidad que una palabra sin relación, como «ventana». Con esas probabilidades, el modelo también evalúa un texto completo: la probabilidad
de una secuencia es el producto de las de cada token dados los anteriores,
\begin{equation}
    p(x_1, \dots, x_n) = \prod_{t=1}^{n} p(x_t \mid x_1, \dots, x_{t-1}),
    \label{eq:marco_lm}
\end{equation}
donde $x_t$ es el token de la posición $t$ y $n$, la longitud de la secuencia~\cite{brown2020language}.

Esas probabilidades se calculan con el decodificador del transformer, que apila los bloques de la
Sección~\ref{sec:marco_atencion} con una diferencia: una máscara en la atención impide que cada token tome
información de los que vienen después. Así, la salida de cada posición depende solo de los tokens anteriores, y
el modelo aprende a predecir el token siguiente en todas las posiciones de un texto a la vez. A diferencia del
aprendizaje supervisado de la Sección~\ref{sec:marco_fundamentos}, esta tarea no necesita que una persona anote
cada ejemplo, porque la etiqueta de cada posición es el token que sigue en el propio texto. Un entrenamiento en el que las etiquetas se obtienen de los mismos datos se llama \textbf{aprendizaje
autosupervisado} (\emph{self-supervised learning}), y permite usar como ejemplo de entrenamiento cualquier texto
disponible~\cite{vaswani2017attention,brown2020language}.

Según la parte del transformer que conservan, los modelos de lenguaje preentrenados son de dos tipos. Los de
solo decodificador se entrenan como se acaba de describir y generan texto. Un \textbf{\gls{llm}} es un modelo de
este tipo preentrenado sobre cientos de miles de millones de tokens, y GPT-3, de 2020, llevó el esquema a 175 mil
millones de parámetros~\cite{brown2020language}. Los de solo codificador, como BERT, conservan la otra mitad:
cada token toma información de los anteriores y de los posteriores, y el modelo se preentrena adivinando tokens
ocultos al azar en medio del texto, también sin anotación manual. Por eso no generan texto, pero representan su
significado, y la Sección~\ref{sec:marco_eval_texto} usa uno para comparar reportes~\cite{devlin2019bert}.

Un \gls{llm} preentrenado, sin embargo, solo continúa textos. Para que responda a pedidos, se le aplica el ajuste fino de la Sección~\ref{sec:marco_transferencia} sobre pares
instrucción--respuesta, y después una alineación con \gls{rlhf}, que lo entrena a preferir las respuestas que
evaluadores humanos marcan como mejores~\cite{ouyang2022training}.

Un modelo que sigue instrucciones se controla con el texto que recibe. Todo lo que se le entrega antes de
generar, la instrucción y cualquier información adicional, se llama \emph{prompt}. Como la probabilidad de la
Ecuación~\ref{eq:marco_lm} está condicionada a los tokens anteriores, cambiar el prompt cambia lo que el modelo
genera sin modificar sus pesos. Los modelos grandes aprovechan además la información y los ejemplos incluidos en
el prompt para resolver una tarea sin ningún paso de entrenamiento, lo que se conoce como aprendizaje en
contexto (\emph{in-context learning})~\cite{brown2020language}. Agregar evidencia al prompt es, por lo tanto, una
forma de condicionar lo que el modelo genera, y la Sección~\ref{sec:marco_rep_generacion} la usa para anclar el
reporte.

Dado un prompt, el texto se genera token a token y en cada paso el modelo devuelve una distribución sobre el
vocabulario, se elige un token y se lo agrega al prompt. La regla con que se elige es la \textbf{estrategia de
decodificación}. La decodificación voraz (\emph{greedy decoding}) toma siempre el token más probable, y en la
generación de texto libre, como los demás métodos que buscan el texto más probable, tiende a repetir
frases~\cite{holtzman2020curious}. Los métodos de muestreo (\emph{sampling}), en cambio, sortean el token
siguiente según la distribución del modelo. La temperatura $T$ divide los logits antes de la softmax de la
Ecuación~\ref{eq:marco_softmax},
\begin{equation}
    p_i = \frac{\exp(z_i / T)}{\sum_{j} \exp(z_j / T)},
    \label{eq:marco_temperatura}
\end{equation}
de modo que con $T<1$ la probabilidad se concentra en los tokens más probables y con $T>1$ se reparte entre más
candidatos. La Sección~\ref{sec:marco_eval_calib} usa la misma operación para calibrar un clasificador. El
muestreo \emph{top-p} sortea solo entre el conjunto más chico de tokens cuya probabilidad acumulada alcanza un
umbral $p$, de modo que la cantidad de candidatos crece cuando la distribución está repartida y se achica cuando
está concentrada. Así evita tanto la repetición de la decodificación voraz como las palabras incoherentes que
aparecen al muestrear de la distribución completa~\cite{holtzman2020curious}. Como la estrategia de
decodificación cambia el texto generado, forma parte de la configuración de un generador.

\subsection{Ajuste eficiente de parámetros: LoRA}
\label{sec:marco_llm_lora}

Para adaptar a la dermatología un modelo de pesos abiertos, es decir, uno cuyos pesos se pueden descargar y
modificar, el ajuste fino de la Sección~\ref{sec:marco_transferencia} necesita más memoria de la que ofrece una
GPU de consumo. Un modelo de 7 mil millones de parámetros ocupa 14~GB solo para almacenar sus pesos con 16 bits, es decir, 2 bytes por parámetro. El ajuste
fino completo agrega un gradiente por parámetro, del mismo tamaño que los pesos, y los dos promedios que Adam
guarda por cada parámetro (Sección~\ref{sec:marco_redes}), de modo que la memoria necesaria llega al menos a
cuatro veces la del modelo, unos 56~GB, sin contar los valores intermedios de la red. Esa cifra es más del
doble de los 24~GB de una GPU de consumo de gama alta. Los métodos de \gls{peft} resuelven este problema:
congelan el modelo preentrenado y entrenan solo un conjunto pequeño de parámetros, de modo que el gradiente y
los promedios de Adam se guardan solo para ellos~\cite{hu2021lora}.

\emph{\Gls{lora}} es uno de esos métodos. Parte de la hipótesis de que el cambio que el ajuste fino introduce en una
matriz de pesos tiene rango bajo, es decir, se puede escribir como el producto de dos matrices delgadas. Sobre
cada matriz congelada $W$, de tamaño $d \times k$, que se elige adaptar, entrena entonces un incremento de bajo
rango,
\begin{equation}
    W' = W + \frac{\alpha}{r}\,BA, \qquad B \in \mathbb{R}^{d\times r},\quad A \in \mathbb{R}^{r\times k},\quad
    r \ll \min(d, k),
    \label{eq:marco_lora}
\end{equation}
y solo actualiza $A$ y $B$. Por ejemplo, para una matriz de $4096 \times 4096$
y $r = 8$, el incremento tiene 65\,536 parámetros, 256 veces menos que la matriz. $A$ se inicializa con valores
aleatorios y $B$ en cero, de modo que el incremento vale cero al empezar y el modelo parte exactamente del
preentrenado. El rango $r$ fija cuántos parámetros tiene el incremento y, con ello, su capacidad de adaptación,
y el factor $\alpha$ escala su efecto, de modo que se puede cambiar $r$ sin volver a ajustar la tasa de
aprendizaje~\cite{hu2021lora}.

% Figura: reparametrización de LoRA. Adaptada de hu2021lora (Figura 1).
% Se usa en sec:marco_llm_lora, después de la Ecuación de LoRA.

\begin{figure}[H]
    \centering
    \begin{tikzpicture}[x=1mm, y=1mm]

        \node[figentrada, text width=10mm, minimum height=7mm] (x) at (0, 0) {$\mathbf{x}$};

        % Rama congelada.
        \node[fignodo, text width=36mm, minimum height=24mm, fill=figrelleno] (W) at (-24, 25)
            {pesos preentrenados\\$W \in \mathbb{R}^{d\times k}$\\[2pt]{\footnotesize congelados}};

        % Rama entrenable: A reduce a r, B vuelve a d.
        \draw[draw=fignaranja!70!black, fill=fignaranja!25, thick] (8, 16) -- (32, 16) -- (26, 22) -- (14, 22) -- cycle;
        \node[font=\footnotesize] at (20, 19) {$A \in \mathbb{R}^{r\times k}$};
        \draw[draw=fignaranja!70!black, fill=fignaranja!25, thick] (14, 28) -- (26, 28) -- (32, 34) -- (8, 34) -- cycle;
        \node[font=\footnotesize] at (20, 31) {$B \in \mathbb{R}^{d\times r}$};
        \node[font=\footnotesize, text=fignaranja!70!black, align=left, anchor=west] at (34, 25)
            {entrenables\\rango $r \ll \min(d,k)$};

        % Suma y salida.
        \node[circle, draw=figtrazo, inner sep=1pt, font=\small] (sum) at (0, 46) {$+$};
        \node[figsalida, text width=26mm, minimum height=7mm] (h) at (0, 58) {$\mathbf{h} = W\mathbf{x} + BA\mathbf{x}$};

        \draw[figflecha] (x.north) -- ++(0, 4) -| (W.south);
        \draw[figflecha] (x.north) -- ++(0, 4) -| (20, 16);
        \draw[figflecha] (20, 22) -- (20, 28);
        \draw[figflecha] (W.north) |- (sum.west);
        \draw[figflecha] (20, 34) |- (sum.east);
        \draw[figflecha] (sum.north) -- (h.south);

    \end{tikzpicture}
    \caption{Reparametrización de LoRA con el incremento de bajo rango $BA$. Adaptada de \textcite{hu2021lora}.}
    \label{fig:marco_lora}
\end{figure}


La Figura~\ref{fig:marco_lora} muestra cómo el incremento $BA$ se suma a la salida de la matriz congelada. El
método no fija qué matrices adaptar. En un transformer, casi todos los parámetros están en las matrices de
proyección de la atención y del perceptrón de cada bloque~\cite{kaplan2020scaling}, y son esas las candidatas.
El trabajo original aplicó \gls{lora} solo a las de la atención (Sección~\ref{sec:marco_atencion}), en la
mayoría de los experimentos a las proyecciones de la consulta y del valor, aunque señala que se puede aplicar a
cualquier subconjunto de matrices. Cuantas más matrices se adaptan y mayor es $r$, más parámetros se entrenan,
pero siguen siendo una fracción mínima del modelo: en GPT-3, \gls{lora} redujo unas 10\,000 veces los
parámetros entrenables y unas tres veces la memoria de la GPU. Al terminar, el incremento se suma a $W$, y el
modelo resultante se ejecuta igual que el original, sin tiempo de cálculo adicional~\cite{hu2021lora}. Un
\gls{llm} adaptado de este modo, sin embargo, solo recibe texto, y un reporte describe una imagen.

\subsection{Modelo visión--lenguaje contrastivo CLIP}
\label{sec:marco_vlm}
\label{sec:marco_vlm_clip}

Para que el modelo reciba también la imagen hace falta un \gls{vlm}, un modelo que procesa imagen y texto. Hay
dos formas de construirlo. La primera entrena un
codificador de imagen y uno de texto para que representen las dos modalidades en un mismo espacio de
embeddings, en el que el embedding de una imagen se puede comparar con el de un texto. El resultado sirve para
clasificar imágenes y para buscar casos parecidos. La segunda conecta un codificador de
imagen a la entrada de un \gls{llm}, que entonces genera texto a partir de la imagen. En ambas, el codificador
de imagen suele ser un \gls{vit} como el de la Sección~\ref{sec:marco_vit}~\cite{radford2021learning,liu2023visual}.

El \gls{clip} es el método de referencia de la primera forma. Entrena un codificador de imagen y uno de texto
sobre pares imagen--texto, cada uno con una proyección final al espacio común. La comparación entre el
embedding de una imagen, $\mathbf{u}$, y el de un texto, $\mathbf{v}$, se mide con la similitud coseno.
\begin{equation}
    \cos(\mathbf{u}, \mathbf{v}) = \frac{\mathbf{u}^{\top}\mathbf{v}}{\lVert\mathbf{u}\rVert\,\lVert\mathbf{v}\rVert},
    \label{eq:marco_cos}
\end{equation}
Esta fórmula es igual a 1 si los vectores apuntan en la misma dirección, 0 si son ortogonales y $-1$ si son opuestos. En cada
lote de $N$ pares, el modelo calcula la similitud entre cada imagen y cada texto, y aprende a distinguir los $N$
pares correctos entre las $N^2$ combinaciones posibles (Figura~\ref{fig:marco_clip}). Para ello, aumenta la
similitud de cada imagen con su propio texto respecto de la que tiene con los textos de los demás pares, y en
ese sentido imagen y texto quedan alineados. El modelo original se entrenó así sobre 400 millones de pares
recogidos de internet, sin ninguna etiqueta de clase~\cite{radford2021learning}.

% Figura: entrenamiento contrastivo de CLIP. Adaptada de radford2021learning (Figura 1, panel 1).
% Se usa en sec:marco_vlm_clip, después del párrafo que define el objetivo contrastivo.

\begin{figure}[H]
    \centering
    \begin{tikzpicture}[x=1mm, y=1mm]

        % Codificador de texto (arriba) y sus embeddings T_1..T_4.
        \node[figentrada, text width=22mm, minimum height=7mm] (txt) at (-26, 44) {textos del lote};
        \node[figtarea, text width=30mm, minimum height=7mm] (enct) at (22, 44) {codificador de texto};
        \draw[figflecha] (txt.east) -- (enct.west);
        \foreach \j in {1,...,4} {
            \node[draw=figverde!70!black, fill=figverde!10, minimum width=9mm, minimum height=6mm,
                font=\footnotesize] (T\j) at (\j*10, 33) {$T_{\j}$};
        }
        \draw[figflecha] (enct.south) -- ++(0, -3.5);

        % Codificador de imagen (izquierda) y sus embeddings I_1..I_4.
        \node[figentrada, text width=18mm, minimum height=7mm] (img) at (-40, 15) {imágenes del lote};
        \node[figtarea, text width=18mm, minimum height=7mm] (enci) at (-17, 15) {codificador de imagen};
        \draw[figflecha] (img.east) -- (enci.west);
        \foreach \i in {1,...,4} {
            \node[draw=figazul!80!black, fill=figazul!12, minimum width=9mm, minimum height=6mm,
                font=\footnotesize] (I\i) at (0, 33-\i*8) {$I_{\i}$};
        }
        \draw[figflecha] (enci.east) -- (I2.west |- enci.east);

        % Matriz de similitudes; la diagonal son los pares correctos.
        \foreach \i in {1,...,4} {
            \foreach \j in {1,...,4} {
                \ifnum\i=\j
                    \node[draw=fignaranja!70!black, fill=fignaranja!30, thick, minimum width=9mm,
                        minimum height=6mm, font=\scriptsize] at (\j*10, 33-\i*8)
                        {$I_{\i}{\cdot}T_{\j}$};
                \else
                    \node[draw=figtrazo, fill=figrelleno, minimum width=9mm, minimum height=6mm,
                        font=\scriptsize, text=figtrazo] at (\j*10, 33-\i*8) {$I_{\i}{\cdot}T_{\j}$};
                \fi
            }
        }
        \node[figetiqueta, align=left, anchor=west] at (47, 17)
            {diagonal: pares\\correctos, se acercan};
        \node[figetiqueta, align=left, anchor=west] at (47, 5)
            {resto: pares\\distintos, se alejan};

    \end{tikzpicture}
    \caption{Entrenamiento contrastivo de CLIP. Adaptada de \textcite{radford2021learning}.}
    \label{fig:marco_clip}
\end{figure}


Ese espacio compartido permite clasificar sin entrenar un clasificador. Cada clase se escribe como un texto, por
ejemplo «una imagen dermatoscópica de un melanoma», y la imagen se asigna a la clase cuyo embedding de texto
tiene mayor similitud coseno con el suyo. Clasificar sin ningún ejemplo etiquetado de la tarea se llama
clasificación \textbf{\emph{zero-shot}}, y su resultado depende de cómo se redacten esos textos: cambiar la plantilla del prompt cambia la
exactitud~\cite{radford2021learning}.

Cuando hay ejemplos etiquetados, el codificador de imagen congelado se puede usar como extractor de
características, con el linear probe de la Sección~\ref{sec:marco_transferencia}, o sin entrenar ningún
parámetro, con el método de los \gls{knn}. Con una base de $m$ casos etiquetados, para una imagen consulta $x$,
este método busca los $k$ casos cuyo embedding $f(x_j)$ tiene mayor similitud coseno con el suyo,
\begin{equation}
    \mathcal{N}_k(x) = \operatorname*{arg\,max}_{S \subseteq \{1,\dots,m\},\ |S| = k}\ \sum_{j \in S}
    \cos\bigl(f(x),\, f(x_j)\bigr),
    \label{eq:marco_vecinos}
\end{equation}
y estima la probabilidad de cada clase $c$ como la proporción de vecinos que la tienen,
\begin{equation}
    \hat{p}_{\text{kNN}}(c \mid x) = \frac{1}{k}\sum_{j \in \mathcal{N}_k(x)} \mathbf{1}[\,y_j = c\,],
    \label{eq:marco_knn}
\end{equation}
donde $y_j$ es la etiqueta del caso $j$. La predicción es la clase de mayor proporción, el voto mayoritario, que
extiende a $k$ vecinos la regla del vecino más cercano~\cite{cover1967nearest}. Así, el espacio contrastivo
sirve también de índice de casos parecidos, que la Sección~\ref{sec:marco_rep_recuperacion} usa como evidencia
para el texto.

Un modelo contrastivo solo relaciona lo que vio en sus pares de entrenamiento, y los pares recogidos de internet
casi no contienen imágenes dermatoscópicas. DermLIP, de 2025, aplica el mismo esquema a alrededor de un millón
de pares imagen--texto dermatológicos, y en clasificación zero-shot de enfermedades de la piel supera a los
modelos biomédicos de propósito general~\cite{yan2025derm1miccv}. Un modelo contrastivo, sin embargo, no genera
texto: compara una imagen con textos dados, pero no puede escribir uno nuevo. Para eso hace falta la segunda
forma de construir un \gls{vlm}.

\subsection{Modelos visión--lenguaje generativos}
\label{sec:marco_vlm_generativos}

LLaVA, de 2023, es un esquema representativo de los \glspl{vlm} generativos abiertos. Un codificador visual del
tipo de \gls{clip} extrae los rasgos de la imagen, y un proyector, en LLaVA una capa lineal, los lleva a la
dimensión de los embeddings de token del \gls{llm}. Los tokens visuales resultantes se concatenan con la
instrucción en el prompt de un \gls{llm} de solo decodificador, que genera la respuesta token a
token. La imagen entra así como una parte más del prompt~\cite{liu2023visual}.

LLaVA se entrena en dos etapas. En la primera, con el codificador visual y el \gls{llm} congelados, solo se
ajusta el proyector sobre pares imagen--leyenda, para que los tokens visuales se parezcan a los embeddings de
palabras que el \gls{llm} ya sabe procesar. En la segunda, el ajuste por instrucciones visuales (\emph{visual
instruction tuning}), se ajustan el proyector y el \gls{llm} sobre conversaciones en las que se pide describir
una imagen, razonar sobre ella o responder preguntas~\cite{liu2023visual}. Como el \gls{llm} concentra la
mayoría de los parámetros, en esa segunda etapa se le puede aplicar \gls{lora}
(Sección~\ref{sec:marco_llm_lora}) en lugar del ajuste completo, mientras el codificador visual sigue congelado
y el proyector se entrena entero.

% Figura: arquitectura de un VLM generativo tipo LLaVA. Adaptada de liu2023visual (Figura 1).
% Se usa en sec:marco_vlm_generativos, después del párrafo que presenta el esquema.
% [H]: como flotante saltaba a la página siguiente y cortaba el párrafo que venía después.

\begin{figure}[H]
    \centering
    \begin{tikzpicture}[x=1mm, y=1mm]

        % Rama visual.
        \node[figentrada, text width=18mm] (img) at (0, 0) {imagen};
        \node[figtarea, text width=22mm] (enc) at (0, 16) {codificador visual\\{\footnotesize (tipo CLIP)}};
        \node[figtarea, text width=22mm] (proj) at (0, 32) {proyector\\{\footnotesize (capa lineal)}};
        \draw[figflecha] (img) -- (enc);
        \draw[figflecha] (enc) -- (proj);

        % Rama de texto.
        \node[figentrada, text width=22mm] (instr) at (40, 0) {instrucción};
        \node[fignodo, text width=22mm, fill=figrelleno] (tok) at (40, 16) {tokenizador y\\\emph{embeddings}};
        \draw[figflecha] (instr) -- (tok);

        % Secuencia de entrada al LLM: cada grupo de tokens queda encima de la rama que lo produce, y un
        % marco agrupa sus casillas para que la flecha llegue al grupo y no a una casilla suelta.
        \foreach \k in {0,...,3} {
            \node[draw=figazul!80!black, fill=figazul!12, minimum width=4.5mm, minimum height=4.5mm]
                (v\k) at (-8.25+\k*5.5, 47) {};
            \node[draw=figtrazo, fill=figrelleno, minimum width=4.5mm, minimum height=4.5mm]
                (t\k) at (31.75+\k*5.5, 47) {};
        }
        \node[draw=figtrazo, densely dotted, rounded corners=1pt, inner sep=1.2mm, fit=(v0) (v3)] (vis) {};
        \node[draw=figtrazo, densely dotted, rounded corners=1pt, inner sep=1.2mm, fit=(t0) (t3)] (txt) {};
        \node[figetiqueta, anchor=east, align=right] at (vis.west) {tokens\\visuales};
        \node[figetiqueta, anchor=west, align=left] at (txt.east) {tokens\\de texto};
        \draw[figflecha] (proj.north) -- (vis.south);
        \draw[figflecha] (tok.north) -- (txt.south);

        % LLM y salida: los dos grupos entran como una sola secuencia.
        \node[figtarea, minimum width=66mm] (llm) at (20, 64) {LLM de solo decodificador};
        \draw[figflecha] (vis.north) -- (vis.north |- llm.south);
        \draw[figflecha] (txt.north) -- (txt.north |- llm.south);
        \node[figsalida, text width=30mm] (out) at (20, 80) {texto generado};
        \draw[figflecha] (llm) -- (out);

        \node[figetiqueta, text=fignaranja!70!black, align=left, anchor=west] at (62, 32)
            {etapa 1: se entrena\\solo el proyector};
        \node[figetiqueta, text=fignaranja!70!black, align=left, anchor=west] at (62, 64)
            {etapa 2: proyector\\y LLM};

    \end{tikzpicture}
    \caption{Arquitectura de un \gls{vlm} generativo del tipo de LLaVA. Adaptada de \textcite{liu2023visual}.}
    \label{fig:marco_llava}
\end{figure}


La Figura~\ref{fig:marco_llava} resume la arquitectura y qué componentes se entrenan en cada etapa. La misma
receta se llevó después al dominio médico cambiando los datos. LLaVA-Med, de 2023, la repite con
figuras y leyendas de artículos biomédicos, y se entrena en menos de quince horas con ocho
GPU~\cite{li2023llavamed}. Su versión 1.5 usa Mistral~7B como modelo de lenguaje~\cite{jiang2023mistral,zeng2025mmskin}. SkinVL,
de 2025, repite la receta con conjuntos dermatológicos, entre ellos uno construido a partir de libros de texto de
la especialidad~\cite{zeng2025mmskin}.

Un \gls{vlm} generativo, en cualquiera de estas variantes, devuelve texto y no una probabilidad por clase. Para
clasificar con él, se puede buscar el nombre de la clase en la respuesta generada o, sin generar, calcular con la
Ecuación~\ref{eq:marco_lm} la probabilidad que el modelo asigna al nombre de cada clase como continuación del
prompt y elegir la mayor. Esa probabilidad se normaliza por la cantidad de tokens del nombre, porque cada token
agrega al producto un factor menor que uno, y sin normalizar se favorecerían los nombres cortos. Los puntajes
que se obtienen ordenan las clases, pero nada garantiza que reflejen la probabilidad de acertar.

El problema no es exclusivo de los \glspl{vlm}. Los clasificadores de la Sección~\ref{sec:marco_vision} también
devuelven una clase para cualquier imagen, incluso cuando se equivocan, y su probabilidad máxima no indica por sí
sola cuánto confiar en esa respuesta~\cite{guo2017calibration}. En el enfoque de este trabajo, resulta necesario establecer cuándo una predicción ofrece evidencia suficiente para utilizarla como diagnóstico en el reporte.
